\section{Datasets}

\subsection{Constellation Diagram Generation}

For modulation classification tasks, constellation diagrams provide superior insights compared to raw time-series data by capturing more detailed signal characteristics \cite{doan2020learning}. The constellation diagram generation process involves several key steps. Initially, modulated signals are projected onto a complex plane, balancing detail preservation with computational efficiency for SNR levels between -10 dB and 10 dB. The visualization process then evolves through multiple phases: beginning with a basic gray-scale representation that handles varying sample densities per pixel, followed by an enhanced version incorporating an exponential decay function ($B_{i,j}$). This enhancement accounts for both the precise sample locations within pixels and their neighborhood effects, considering factors like sample power ($P$), sample-to-pixel centroid distances ($d_{i,j}$), and decay rate ($\alpha$) \cite{peng2018modulation}. To align with DenoMAE2.0's input requirements, we create a three-channel RGB image by generating three separate enhanced grayscale images from identical samples, each using different decay parameters.

Our dataset includes ten modulation types as outlined in Table~\ref{tab:modulation_classes}. The pretraining dataset comprises 10,000 samples uniformly distributed across ten modulation types, with randomly assigned SNR values in the range of -10 dB to 10 dB. For the downstream classification task, we use 1000 samples for training and 100 samples for testing, evenly distributed among the 10 classes, evaluated at SNR values from -10 dB to 10 dB in 1 dB increments.

\begin{table}[htbp]
    \centering
    \caption{List of all modulation classes}
    \begin{tabularx}{\linewidth}{Xl}
    \hline
    \textbf{Modulation Technique} & \textbf{Abbreviation} \\
    \hline
    4-Amplitude Shift Keying & 4ASK \\
    4-Pulse Amplitude Modulation & 4PAM \\
    8-Amplitude Shift Keying & 8ASK \\
    16-Pulse Amplitude Modulation & 16PAM \\
    Continuous-Phase Frequency Shift Keying & CPFSK \\
    Differential Quadrature Phase Shift Keying & DQPSK \\
    Gaussian Frequency Shift Keying & GFSK \\
    Gaussian Minimum Shift Keying & GMSK \\
    On-Off Keying & OOK \\
    Offset Quadrature Phase Shift Keying & OQPSK \\
    \hline
    \end{tabularx}
    \label{tab:modulation_classes}
\end{table}

% Our dataset includes ten modulation types:Four-Level Amplitude Shift Keying (4ASK), Four-Level Pulse Amplitude Modulation (4PAM), Eight-Level Amplitude Shift Keying (8ASK), Sixteen-Level Pulse Amplitude Modulation (16PAM), Continuous Phase Frequency Shift Keying (CPFSK), Differential Quadrature Phase Shift Keying (DQPSK), Gaussian Frequency Shift Keying (GFSK), Gaussian Minimum Shift Keying (GMSK), On-Off Keying (OOK), and Offset Quadrature Phase Shift Keying (OQPSK). 

\subsection{Parameters in Sample Generation}

Our dataset incorporates ten distinct modulation formats. Each signal initially has a length of $L_0 = 1024$ samples. The preprocessing pipeline adapts these signals for transformer architecture compatibility through two main steps: first transforming into $\mathbf{S}_1 \in \mathbb{R}^{32 \times 32}$, then interpolating to $\mathbf{S}_2 \in \mathbb{R}^{224 \times 224}$. For modulation formats exceeding $L_0$ (such as GMSK with $L_\text{GMSK} = 8196$), we perform initial downsampling to $L_0$. The final signal representation $\mathbf{S}_3 \in \mathbb{R}^{3 \times 224 \times 224}$ is created by triplicating $\mathbf{S}_2$ to match the three-channel structure of constellation images $\mathbf{I} \in \mathbb{R}^{3 \times 224 \times 224}$. A uniform sampling rate of $f_s = 200$ kHz is maintained across all modulation formats.

\subsection{Model Parameters and Hyperparameters}

DenoMAE2.0 architecture uses distinct configurations for pretraining and fine-tuning phases (as shown in Table \ref{tab:pre_para} and Table \ref{tab:fine_para}). The model processes $224 \times 224$ pixel images divided into $16 \times 16$ patches. During pretraining, we employ an encoder-decoder structure where the encoder features 12 layers with 12 attention heads and an embedding dimension of 768, while the decoder is lighter with 8 layers, 8 attention heads, and a 512-dimensional embedding space. A mask ratio of 0.75 is applied during the self-supervised pretraining phase, balancing reconstruction difficulty with information preservation.

For fine-tuning, only the encoder portion is retained with its 12-layer, 12-attention-head configuration maintained, adapting the output for either 10-class (original dataset) or 24-class (extended dataset) classification tasks. The learning rate is reduced from 0.0003 during pretraining to 0.0001 for fine-tuning, reflecting the transition from broad feature learning to targeted classification optimization. Both phases maintain a weight decay of 0.05 for regularization, but batch sizes differ (64 for pretraining vs. 32 for fine-tuning) to accommodate the computational demands of each training regime.

\begin{table}[htbp]
    \centering
    \caption{List of hyperparameters during pretraining}
    \label{tab:pre_para}
    \begin{tabular}{cc}
        \hline
        \textbf{Parameter name} & \textbf{Value} \\ \hline
        Training batch size & 64 \\ 
        Testing batch size & 10 \\ 
        Image size & 224, 224 \\ 
        Patch size & 16 \\ 
        Mask ratio & 0.75 \\ 
        Encoder embedding dimension & 768 \\ 
        Decoder embedding dimension & 512 \\ 
        Encoder depth & 12 \\ 
        Decoder depth & 8 \\ 
        Encoder number of heads & 12 \\ 
        Decoder number of heads & 8 \\ 
        Number of epochs & 100 \\ 
        Learning rate & 0.0003 \\ 
        Weight decay & 0.05 \\ 
        Classification loss weight & 0.1 \\ 
        Reconstruction loss weight & 1.0 \\ 
        \hline
    \end{tabular}
\end{table}

\begin{table}[htbp]
    \centering
    \caption{List of hyperparameters during fine-tuning}
    \label{tab:fine_para}
    \begin{tabular}{cc}
        \hline
        \textbf{Parameter name} & \textbf{Value} \\ \hline
        Training batch size & 32 \\ 
        Testing batch size & 10 \\
        Image size & 224, 224 \\ 
        Patch size & 16 \\ 
        Encoder embedding dimension & 768 \\ 
        Encoder depth & 12 \\ 
        Encoder number of heads & 12 \\ 
        Number of epochs & 150 \\
        Learning rate & 0.0001 \\ 
        Weight decay & 0.05 \\ 
        Number of classes & 10 \\ \hline
    \end{tabular}
\end{table}


